\subsection{Модуль Ledbar}

Для управления светодиодами в системе Geoskan Pioneer используется встроенный модуль \texttt{Ledbar}. Он предоставляет простой интерфейс (API) для работы с цветом.

\begin{codebox}[Основные функции]
\begin{itemize}
  \item \texttt{Ledbar.new(N)} — конструктор. Создаёт объект для управления лентой из \texttt{N} светодиодов.
  \begin{itemize}
      \item \textit{Пример:} \texttt{local leds = Ledbar.new(25)}
  \end{itemize}
  
  \item \texttt{leds:set(i, r, g, b)} — устанавливает цвет конкретного светодиода.
  \begin{itemize}
      \item \texttt{i} — индекс светодиода (от \texttt{0} до \texttt{N-1}).
      \item \texttt{r, g, b} — яркость каналов в диапазоне от \texttt{0.0} (выключен) до \texttt{1.0} (максимум).
  \end{itemize}
\end{itemize}
\end{codebox}

\begin{warnbox}[Индексация]
Обратите внимание: нумерация светодиодов начинается с \textbf{нуля}. Если вы создали ленту на 25 элементов, допустимые индексы — от 0 до 24. Попытка обратиться к индексу 25 приведёт к ошибке.
\end{warnbox}

Также часто используется вспомогательная функция для заливки всей ленты одним цветом:
\begin{minted}{lua}
local function changeColor(color)
  for i = 0, ledNumber - 1 do
    leds:set(i, table.unpack(color))
  end
end
\end{minted}
